\documentclass[a4paper, serif, 10pt]{moderncv}

%% moderncv
% color and style
\moderncvcolor{black}
\moderncvstyle{banking}

%% page layout/margins
\usepackage[top=2.5cm, bottom=2.5cm, left=1.5cm, right=1.5cm]{geometry}

\nopagenumbers{}

%% Babel config for Dutch language
% ref:
% - https://latex3.github.io/babel/guides/locale-dutch.html
\usepackage[dutch]{babel}

%% fontspec
\usepackage{fontspec}

\IfFontExistsTF{Linux Libertine}{
  \setmainfont[Ligatures={TeX, Common}, Numbers=OldStyle]{Linux Libertine}
}{}
\IfFontExistsTF{Linux Libertine O}{
  \setmainfont[Ligatures={TeX, Common}, Numbers=OldStyle]{Linux Libertine O}
}{}

%% CTeX
% CJK support, used to show CJK characters in the resume
%
% - heading=false: disable ctex heading
% - punct=kaiming: use kaiming punctuations styles for CJK
% - scheme=plain: use plain scheme, do not override \normalsize font size
% - space=auto: space settings for CJK characters
%
% CJK fonts are resolved per locale language, e.g. Noto Serif/Sans CJK JP
% for Japanese resumes, with fallback to the next region if not installed
%
% ref:
% - http://ctan.mirrorcatalogs.com/language/chinese/ctex/ctex.pdf
\usepackage[UTF8, heading=false, punct=kaiming, scheme=plain, space=auto]{ctex}

\IfFontExistsTF{Noto Serif CJK SC}{
  \setCJKmainfont{Noto Serif CJK SC}
}{}
\IfFontExistsTF{Noto Sans CJK SC}{
  \setCJKsansfont{Noto Sans CJK SC}
}{}

%% line spacing
\usepackage{setspace}
\setstretch{1.125}

%% URL styling - use normal text instead of monospace
\urlstyle{same}

%% Patch \httplink and \httpslink to handle full URLs
% The original moderncv macros always prepend http:// or https:// to URLs.
% This causes issues when the URL already has a protocol (e.g., https://example.com)
% resulting in malformed URLs like https://https://example.com.
% This patch checks if the URL already contains "://" and uses it directly if so.
% Note: We use \str_set:Nx (x-expansion) to expand macros like \@homepage first.
\ExplSyntaxOn
\str_new:N \g_yamlresume_url_str

\RenewDocumentCommand{\httpslink}{O{}m}{
  \str_set:Nx \g_yamlresume_url_str {#2}
  \str_if_in:NnTF \g_yamlresume_url_str {://}
    { \url{#2} }
    { \url{https://#2} }
}

\RenewDocumentCommand{\httplink}{O{}m}{
  \str_set:Nx \g_yamlresume_url_str {#2}
  \str_if_in:NnTF \g_yamlresume_url_str {://}
    { \url{#2} }
    { \url{http://#2} }
}
\ExplSyntaxOff

\name{Sanne de Vries}{}
\title{Operations Manager met focus op logistiek en procesverbetering}
\phone[mobile]{+31 6 1234 5678}
\email{sanne.devries@example.nl}
\homepage{https://www.sannedevries.nl}

\address{Keizersgracht 123, Amsterdam, Noord-Holland, Nederland, 1015 CJ}{}{}

\extrainfo{{\small \faLinkedin}\ \href{https://www.linkedin.com/in/sannedevries}{@sannedevries}}

\begin{document}

\maketitle

\section{Basis Informatie}

\cvline{}{Operations Manager met meer dan acht jaar ervaring in logistiek, productie en distributie. Ik combineer analytisch inzicht met daadkrachtig leiderschap en weet operationele prestaties duurzaam te verbeteren.
%
\begin{itemize}
\item Stuurt multidisciplinaire teams aan van 10 tot 50 medewerkers en creëert een cultuur van eigenaarschap en continue verbetering.
\item Realiseert aantoonbare resultaten op het gebied van doorlooptijd, kwaliteit, kosten en klanttevredenheid.
\item Werkt datagedreven met KPI's, Lean-principes en moderne plannings- en magazijnsystemen.
\item Communiceert helder met stakeholders van werkvloer tot directie en vertaalt strategie naar uitvoerbare plannen.
\end{itemize}}
\section{Opleidingen}

\cventry{sep 2010–jul 2014}
        {Bachelor, Logistiek en Economie, Score: 8.1}
        {Hogeschool van Amsterdam}
        {\url{https://www.hva.nl}}
        {}
        {\begin{itemize}
\item Opgedaan in operationeel management, logistieke processen en financiële sturing.
\item Afstudeeronderzoek naar voorraadoptimalisatie bij een regionale distributeur.
\end{itemize}
\textbf{Cursussen}: Supply Chain Management, Operationeel Management, Statistiek voor Logistiek, Projectmanagement, Financieel Management, Transporteconomie}

\cventry{sep 2014–jul 2016}
        {Master, Transport, Infrastructuur en Logistiek, Score: 8.4}
        {Technische Universiteit Delft}
        {\url{https://www.tudelft.nl}}
        {}
        {\begin{itemize}
\item Verdieping in transportnetwerken, simulatie en ketenregie.
\item Scriptie over datagedreven capaciteitsplanning in de foodsector.
\end{itemize}
\textbf{Cursussen}: Transportnetwerken, Simulatie van Logistieke Systemen, Ketenregie, Operationeel Onderzoek, Duurzame Mobiliteit}
\section{Werkervaring}

\cventry{jan 2021–Heden}
        {Operations Manager}
        {FreshFlow Distributie}
        {\url{https://www.freshflowdistributie.nl}}
        {}
        {\begin{itemize}
\item Leidt een operationeel team van 45 medewerkers verdeeld over planning, magazijn en transport.
\item Verantwoordelijk voor dagelijkse aansturing, veiligheid, kwaliteit en leverbetrouwbaarheid.
\item Realiseerde een daling van 18\% in orderdoorlooptijd door herinrichting van pickroutes en planning.
\item Stuurde op KPI's zoals OTIF, productiviteit, voorraadnauwkeurigheid en medewerkerstevredenheid.
\item Beheerde een operationeel budget van EUR 6,5 miljoen en realiseerde 7\% kostenbesparing.
\end{itemize}
\textbf{Trefwoorden}: Operationele excellentie, Teamontwikkeling, KPI-sturing, Budgetverantwoordelijkheid}

\cventry{mrt 2017–dec 2020}
        {Teamleider Logistiek}
        {Nexus Logistiek}
        {\url{https://www.nexuslogistiek.nl}}
        {}
        {\begin{itemize}
\item Gaf leiding aan drie teamleiders en 30 logistieke medewerkers in een 24/7 distributiecentrum.
\item Introduceerde dagelijkse huddle- en visual management-rituelen voor snellere bijsturing.
\item Verhoogde de forecastnauwkeurigheid met 12\% door samenwerking met sales en planning.
\item Begeleidde een verhuizing naar een nieuw magazijn zonder verstoring van klantleveringen.
\end{itemize}
\textbf{Trefwoorden}: Procesverbetering, Verandermanagement, Leveranciersmanagement, Capaciteitsplanning}

\cventry{sep 2016–feb 2017}
        {Procesanalist Operations}
        {Bureau Proceskracht}
        {\url{https://www.bureauproceskracht.nl}}
        {}
        {\begin{itemize}
\item Analyseerde order-, voorraad- en transportprocessen voor middelgrote productie- en retailklanten.
\item Vertaalde data naar concrete verbeterplannen en begeleidde teams bij de implementatie.
\item Ontwikkelde KPI-rapportages en standaardwerkwijzen voor operations en planning.
\end{itemize}
\textbf{Trefwoorden}: Data-analyse, Procesoptimalisatie, Rapportages, Standaardisatie}
\section{Talen}

\cvline{Nederlands}{Moedertaal of tweetalige vaardigheid \hfill \textbf{Trefwoorden}: Moedertaal}
\cvline{Engels}{Volledige professionele vaardigheid \hfill \textbf{Trefwoorden}: C1, IELTS 7.5}
\cvline{Duits}{Beperkte werkvaardigheid \hfill \textbf{Trefwoorden}: B1}
\section{Vaardigheden}

\cvline{Operationeel Management}{Expert \hfill \textbf{Trefwoorden}: Procesoptimalisatie, KPI-dashboards, Capaciteitsplanning, Budgetbeheer}
\cvline{Supply Chain}{Geavanceerd \hfill \textbf{Trefwoorden}: Voorraadbeheer, Vraagvoorspelling, Leveranciersmanagement, Magazijnbeheersysteem}
\cvline{Lean \& Six Sigma}{Geavanceerd \hfill \textbf{Trefwoorden}: Waardestroomanalyse, Kaizen, Oorzaakanalyse, DMAIC}
\cvline{Data \& Rapportage}{Geavanceerd \hfill \textbf{Trefwoorden}: Excel, Power BI, SQL, Tableau}
\cvline{Leiderschap}{Expert \hfill \textbf{Trefwoorden}: Coaching, Verandermanagement, Stakeholdermanagement, Teamontwikkeling}
\section{Onderscheidingen}

\cventry{nov 2022}
        {FreshFlow Distributie}
        {Lean Six Sigma Green Belt Project Award}
        {}
        {}
        {\begin{itemize}
\item Erkend voor een project dat de orderdoorlooptijd met 18\% verlaagde.
\item Beoordeeld op meetbare impact, duurzaamheid van de oplossing en betrokkenheid van het team.
\end{itemize}}

\cventry{dec 2019}
        {Nexus Logistiek}
        {Beste Teamprestatie van het Jaar}
        {}
        {}
        {\begin{itemize}
\item Toegekend aan het logistieke team voor uitstekende leverbetrouwbaarheid tijdens piekperiodes.
\item Gerealiseerd door strakke planning, heldere communicatie en sterke samenwerking.
\end{itemize}}
\section{Certificaten}

\cventry{jun 2020}
        {ASQ}
        {Lean Six Sigma Green Belt}
        {\url{https://asq.org/cert/lean-six-sigma-green-belt}}
        {}
        {}

\cventry{feb 2019}
        {AXELOS}
        {PRINCE2 Foundation}
        {\url{https://www.axelos.com/certifications/prince2}}
        {}
        {}

\cventry{okt 2021}
        {ASCM}
        {Certified Supply Chain Professional (CSCP)}
        {\url{https://www.ascm.org/credentials/cscp/}}
        {}
        {}
\section{Publicaties}

\cventry{mrt 2023}
        {Van voorraadkosten naar klantwaarde: datagedreven operations in de foodsector}
        {Logistiek Magazine}
        {\url{https://www.logistiekmagazine.nl/artikelen/datagedreven-operations}}
        {}
        {\begin{itemize}
\item Praktijkartikel over het verbinden van voorraadsturing aan klantwaarde en servicegraden.
\item Gaat in op KPI-keuzes, forecastverbetering en samenwerking tussen operations en commercie.
\end{itemize}}
\section{Projecten}

\cventry{mrt 2022–nov 2022}
        {Uitrol van een warehouse management systeem voor magazijn, planning en transport.}
        {Implementatie nieuw WMS}
        {\url{https://www.freshflowdistributie.nl/projecten/wms}}
        {}
        {\begin{itemize}
\item Leidde de operationele implementatie van een nieuw WMS binnen een live omgeving.
\item Betrok teamleiders en key users bij procesontwerp, testen en training.
\item Realiseerde een hogere voorraadnauwkeurigheid en kortere inwerktijd voor nieuwe medewerkers.
\end{itemize}
\textbf{Trefwoorden}: WMS, Implementatie, Procesontwerp, Training}

\cventry{jan 2023–jun 2023}
        {Herinrichting van pickroutes en pickmethoden om doorlooptijd en kwaliteit te verbeteren.}
        {Procesoptimalisatie orderpicking}
        {\url{https://www.freshflowdistributie.nl/projecten/orderpicking}}
        {}
        {\begin{itemize}
\item Analyseerde pickdata en observeerde werkprocessen op de werkvloer.
\item Testte nieuwe routeplanning en artikelplaatsing in een pilot.
\item Verlaagde de gemiddelde picktijd met 14\% en verminderde orderfouten.
\end{itemize}
\textbf{Trefwoorden}: Orderpicking, Lean, Data-analyse, Kwaliteit}
\section{Interesses}

\cvline{Duurzaamheid}{Circulaire logistiek, Energiebesparing}
\cvline{Sport}{Hardlopen, Roeien, Wielrennen}
\cvline{Maatschappelijke impact}{Voedselbank, Kennisoverdracht}
\section{Vrijwilligerswerk}

\cventry{feb 2018–Heden}
        {Logistiek Coördinator}
        {Voedselbank Amsterdam}
        {\url{https://www.voedselbankamsterdam.nl}}
        {}
        {\begin{itemize}
\item Coördineert wekelijks de aanvoer, sortering en verdeling van voedingsmiddelen.
\item Stemt planning af met vrijwilligers, winkels en partnerorganisaties.
\item Verbeterde de sorteerprocessen en zorgde voor betrouwbare leveringen aan lokale uitgiftepunten.
\end{itemize}}
    

\cventry{sep 2019–jun 2021}
        {Projectvrijwilliger}
        {Stichting Duurzame Stad}
        {\url{https://www.duurzamestad.nl}}
        {}
        {\begin{itemize}
\item Ondersteunde een project voor efficiëntere stadsdistributie en minder lege kilometers.
\item Verzamelde data, maakte analyses en hielp bij de vertaling naar aanbevelingen voor de gemeente.
\end{itemize}}
    
\end{document}